\begin{salyga}
Stačiojo trikampio $ABC$ statinių ilgiai yra $AC=6\,\text{cm}$ ir
$BC=8\,\text{cm}$, o kampas $\angle ACB$ yra statusis. Iš viršūnės $C$
nuleista aukštinė $CH$ į įžambinę $AB$.
\begin{enumerate}[label=\alph*)]
  \item Apskaičiuokite įžambinės $AB$ ilgį.
  \item Apskaičiuokite trikampio $ABC$ plotą ir aukštinės $CH$ ilgį.
\end{enumerate}
\end{salyga}

\begin{atsakymas}
\begin{enumerate}[label=\alph*)]
  \item $AB=10\,\text{cm}$.
  \item $S_{ABC}=24\,\text{cm}^2$, $CH=4{,}8\,\text{cm}$.
\end{enumerate}
\end{atsakymas}

\begin{sprendimas}
\begin{enumerate}[label=\alph*)]
  \item Trikampis $ABC$ yra statusis, todėl pagal Pitagoro teoremą
  \[
    AB^2=AC^2+BC^2=6^2+8^2=36+64=100.
  \]
  Kadangi kraštinės ilgis yra teigiamas, $AB=10\,\text{cm}$.

  \item Statiniai yra statmeni, todėl vieną statinį laikydami pagrindu,
  o kitą – aukštine, gauname
  \[
    S_{ABC}=\frac{AC\cdot BC}{2}
    =\frac{6\cdot8}{2}=24\,\text{cm}^2.
  \]
  Tą patį plotą išreikškime pasirinkę įžambinę $AB$ pagrindu. Kadangi
  $CH\perp AB$, atkarpa $CH$ yra šio pagrindo aukštinė, taigi
  \[
    S_{ABC}=\frac{AB\cdot CH}{2}.
  \]
  Įstatę apskaičiuotas reikšmes ir išsprendę lygtį gauname
  \[
    24=\frac{10\cdot CH}{2},\qquad
    CH=\frac{2\cdot24}{10}=4{,}8\,\text{cm}.
  \]
  Vadinasi, trikampio plotas yra $24\,\text{cm}^2$, o aukštinės ilgis
  yra $4{,}8\,\text{cm}$.
\end{enumerate}
\hfill$\square$
\end{sprendimas}
